% Source: https://artofproblemsolving.com/community/c3h2572089_after_sieve_of_eranthoses_is_done
% Stable identity: aops-topic-2572089-q01
\textbf{(a)} If $n\le100$ is composite, write $n=pq$ with $1<p\le q$. Then $p^2\le pq=n\le100$, so $p\le10$, and a prime factor of $p$ is one of $2,3,5,7$. So every composite number up to $100$ is a multiple of $2$, $3$, $5$ or $7$ and gets crossed out.

\textbf{(b)} Count the numbers from $1$ to $100$ divisible by none of $2,3,5,7$ by inclusion--exclusion:
\begin{align*}
&100-\left(50+33+20+14\right)+\left(16+10+7+6+4+2\right)-\left(3+2+1+0\right)+0\\
&\qquad=100-117+45-6=22.
\end{align*}
(For example $16=\left\lfloor\dfrac{100}6\right\rfloor$ counts the multiples of $2\cdot3$.) These $22$ numbers are $1$ and the primes from $11$ to $97$. Adding back $2,3,5,7$ and removing $1$, the number of primes up to $100$ is $22-1+4=\boxed{25}$.
